\section{Discussion}

\label{sec:discussion}

\subsection{Implications for Conservation Education}

The results demonstrate that adaptive learning systems with ecological grounding can substantially improve conservation education outcomes in diverse, resource-constrained settings.
Three findings merit particular attention.

First, the 34\% retention advantage of ZooEd-Full over static curricula at 12 months represents a practically significant improvement for organizations investing in long-term workforce development.
Second, the transfer results show that improved retention translates to measurably better field performance, addressing a long-standing critique that educational technology improves test scores without practical impact \citep{stern2014value}.
Third, the ecological grounding mechanism demonstrates that connecting educational content to local environmental events through IoT sensor networks creates powerful contextual learning cues that standard spaced repetition cannot achieve.

\subsection{Scalability Considerations}

The edge node architecture enables ZooEd to operate in environments with intermittent or absent internet connectivity.
Solar-powered edge nodes maintained 99.2\% uptime across deployment sites, with content synchronization occurring during periodic connectivity windows.
The offline-first mobile client cached up to 500 learning activities, sufficient for 2--3 weeks of continuous use without synchronization.

\subsection{Limitations}

Several limitations qualify our findings.
The quasi-experimental design, while strengthened by stratified matching, cannot fully control for self-selection effects.
Assessment reliability for low-resource languages, though improved by calibration, remains below the $\kappa = 0.80$ threshold conventionally required for high-stakes decisions.
The ecological grounding component depends on Zoo sensor network coverage, which was not uniform across sites; three sites had only basic weather station data, limiting the diversity of ecological triggers.

\subsection{Ethical Considerations}

Deploying AI-driven education in indigenous and marginalized communities raises important ethical questions.
We addressed these through participatory design workshops at each site, community veto power over content, on-chain consent with right to deletion, and revenue sharing through token incentives.
However, we acknowledge that the power dynamics inherent in technology deployment cannot be fully resolved through technical mechanisms alone \citep{couldry2019data}.

% ============================================================
% 10. Conclusion
% ============================================================
